\chapter{Conclusion}
\begingroup\small

The empty \emph{Fraud Detection} box is both a warning and an opportunity. It is a warning because a label in an architecture can conceal the absence of purpose, requirements, evidence, accountability, and remedy. It is an opportunity because the missing capability is broader and more valuable than a binary fraud classifier.

This thesis proposes the \product: a federated, purpose-bound layer connecting existing biometric, credential, registry, wallet, device, process, and human evidence. Its scientific core combines supervised known-fraud detection, open-set anomaly detection, relational graph risk, and process-exclusion detection. Its engineering core is calibrated action selection under fraud, friction, human-capacity, privacy, and exclusion constraints. Its governance core separates anomaly from guilt, prevents a universal citizen score, and makes appeal and record repair executable.

The work preserves the essential contributions of \emph{Beyond the Matcher}, the Biometric Fraud Simulation Framework, the Biometric Evaluation Platform, and the National Identity Fraud Risk Engine GNN. The new contribution is their integration into an ecosystem-scale trust and inclusion doctrine that includes credentials such as diplomas and professional qualifications, private and public relying parties, and citizens as rights-bearing participants rather than passive data subjects.

The synthetic demonstrator proves executability, not field performance. Its high metrics reflect planted patterns and must not be marketed as operational evidence. The next scientific step is leakage-resistant retrospective evaluation with real governance and outcome labels. The next industrial step is a narrow observability and shadow-mode pilot on one defined identity or credential journey.

The market already contains strong fraud, biometric, identity-proofing, behavioural, and graph products. A credible product cannot claim an empty competitive field. Its defensible position is the vendor-neutral, sovereign, fingerprint-capable, credential-aware assurance layer that treats exclusion, contestation, and downstream repair as first-class product functions.

The final proposition is therefore:

\begin{quote}
Identity infrastructure is trustworthy only when it resists impersonation and fabricated credentials, detects when its own processes exclude legitimate people, requests no more evidence than necessary, and can repair the consequences of being wrong.
\end{quote}

That proposition is technically challenging, industrially differentiating, and socially worth testing---provided that the work advances through explicit claims, controlled evidence, and the continuing possibility of redesign or no-go.
\endgroup
